\section{Conclusion}

This paper presented a comprehensive comparative analysis of four major machine translation tools (Google Translate, DeepL, Microsoft Translator, Amazon Translate) on technical documentation and UI texts for the English-Turkish language pair.

\subsection{Summary of Findings}

Our evaluation of 10,000 test segments using multiple automatic metrics (BLEU, METEOR, TER, chrF++) revealed:

\begin{enumerate}
    \item \textbf{Tool-Specific Strengths}: [Tool X] excels at technical documentation (BLEU=X.XX), while [Tool Y] performs best on UI strings (BLEU=Y.YY)
    
    \item \textbf{Statistical Significance}: Performance differences between tools are statistically significant (p<0.001), justifying careful tool selection
    
    \item \textbf{Category Matters}: Text category significantly impacts translation quality, with UI strings proving more challenging than technical documentation
    
    \item \textbf{Cost-Quality Trade-offs}: Microsoft Translator offers best cost efficiency (\$10/1M chars), while [Tool X] provides highest quality at premium cost
    
    \item \textbf{Placeholder Handling}: Critical for software functionality; [Tool X] achieves XX\% preservation rate
\end{enumerate}

\subsection{Practical Contributions}

This work provides actionable insights for software development teams:

\begin{itemize}
    \item Evidence-based tool selection guidelines
    \item Open-source dataset for Turkish software localization
    \item Web application for interactive tool comparison
    \item Benchmark scores for quality assessment
\end{itemize}

\subsection{Limitations and Future Work}

While our study provides valuable insights, several directions warrant further investigation:

\subsubsection{Extended Language Pairs}

Future work should examine:
\begin{itemize}
    \item Other Turkic languages (Azerbaijani, Kazakh, Uzbek)
    \item Morphologically rich languages (Finnish, Hungarian)
    \item Low-resource language pairs
\end{itemize}

\subsubsection{Human Evaluation}

Complementing automatic metrics with human evaluation would provide:
\begin{itemize}
    \item Fluency and adequacy ratings
    \item Preference judgments
    \item Error severity classification
    \item Cultural appropriateness assessment
\end{itemize}

\subsubsection{Domain Adaptation}

Investigating fine-tuned models trained on software localization data could:
\begin{itemize}
    \item Improve technical terminology consistency
    \item Enhance placeholder preservation
    \item Better handle domain-specific constructions
\end{itemize}

\subsubsection{Context-Aware Translation}

Developing methods to provide surrounding context for UI strings could address the context dependency challenge identified in our study.

\subsubsection{Temporal Analysis}

Longitudinal studies tracking MT system improvements over time would reveal:
\begin{itemize}
    \item Rate of quality improvement
    \item Stability of tool rankings
    \item Impact of model updates
\end{itemize}

\subsection{Broader Impact}

This research contributes to the growing body of work on MT evaluation for specialized domains. By focusing on software localization, we address the practical needs of a large developer community while advancing understanding of MT system capabilities and limitations.

The open-source release of our dataset, evaluation framework, and web application enables:
\begin{itemize}
    \item Reproducible research
    \item Community-driven improvements
    \item Practical tool adoption
    \item Future comparative studies
\end{itemize}

\subsection{Final Remarks}

As machine translation systems continue to evolve, regular evaluation on domain-specific tasks becomes increasingly important. Our study demonstrates that tool selection should be informed by specific use cases, text characteristics, and project constraints rather than general reputation or popularity.

We hope this work serves as a foundation for future research on MT for software localization and provides practical guidance for developers navigating the complex landscape of translation tools.

\section*{Acknowledgments}

We thank the open-source communities behind OPUS, Mozilla, GNOME, and KDE for providing high-quality localization data. We also acknowledge the developers of SacreBLEU and NLTK for their evaluation metric implementations.

\section*{Data Availability}

The dataset, code, and supplementary materials are available at: [GitHub URL]
